\chapter*{Global Exercise Bank --- Multiple Correct MCQs}
\addcontentsline{toc}{chapter}{Global Exercise Bank: Multiple Correct MCQs}
\label{chap:multiple_correct}

\begin{tcolorbox}[enhanced,colback=subtleblue,colframe=primarynavy,arc=2mm,boxrule=1pt,
    title=\textbf{\large One or More Than One Correct Options (All Chapters)}]
Sources: \textbf{[NRJ]} Ex-II Multi, \textbf{[PRSN]} Ex-II, \textbf{[CNG-TH]}, \textbf{[NA]} Level-3.
\end{tcolorbox}

\begin{problembox}
\textbf{MC.1} \hfill \textbf{[NRJ Ex-II Q5 / PRSN Ex-II Q3]}

Which of the following statements about the molar conductance of electrolytes are \textbf{correct}?
\begin{enumerate}[label=(\alph*)]
    \item For strong electrolytes, $\Lambda_m$ increases linearly with $1/\sqrt{c}$ according to Kohlrausch's law.
    \item For weak electrolytes, $\Lambda_m$ cannot be obtained by extrapolating the $\Lambda_m$ vs $\sqrt{c}$ plot.
    \item The specific conductance ($\kappa$) always decreases on dilution.
    \item At the same concentration, $\Lambda_m$ of a strong electrolyte is always greater than that of a weak electrolyte of the same formula type.
\end{enumerate}
\end{problembox}
\begin{solution}
\textbf{Correct: (a), (b), (c)}

(a) True: Debye--H\"{u}ckel--Onsager equation $\Lambda_m = \Lambda_m^\circ - b\sqrt{c}$ is linear in $\sqrt{c}$.\\
(b) True: Weak electrolytes show a steep non-linear rise at low $c$; cannot be extrapolated.\\
(c) True: $\kappa = \Lambda_m \cdot c/1000$. Though $\Lambda_m$ increases on dilution, $c$ decreases much faster, so $\kappa$ always decreases on dilution.\\
(d) Not always true: At very low $c$, the weak electrolyte may be nearly fully dissociated and approach $\Lambda_m^\circ$, which can be comparable to the strong electrolyte.
\end{solution}

\begin{problembox}
\textbf{MC.2} \hfill \textbf{[PRSN Ex-II Q7 / NRJ Multi Q8]}

Which of the following cells have $E^\circ_{\text{cell}} > 0$ (spontaneous under standard conditions)?
\begin{enumerate}[label=(\alph*)]
    \item $\ce{Zn | Zn^2+ || Cu^2+ | Cu}$; $E^\circ(\ce{Zn^2+/Zn})=-0.76$, $E^\circ(\ce{Cu^2+/Cu})=+0.34\,\text{V}$
    \item $\ce{Cu | Cu^2+ || Ag+ | Ag}$; $E^\circ(\ce{Cu^2+/Cu})=+0.34$, $E^\circ(\ce{Ag+/Ag})=+0.80\,\text{V}$
    \item $\ce{Ag | Ag+ || Cu^2+ | Cu}$ (reverse of above)
    \item $\ce{Pt, H2 | H+(pH=7) | O2(1 bar), Pt}$; $E^\circ(\ce{O2/H2O})=+1.23\,\text{V}$
\end{enumerate}
\end{problembox}
\begin{solution}
\textbf{Correct: (a), (b), (d)}

(a) $E^\circ = 0.34 - (-0.76) = +1.10\,\text{V} > 0$ \checkmark\\
(b) $E^\circ = 0.80 - 0.34 = +0.46\,\text{V} > 0$ \checkmark\\
(c) $E^\circ = 0.34 - 0.80 = -0.46\,\text{V} < 0$ (non-spontaneous) \texttimes\\
(d) At pH = 7: $E_{\ce{H+/H2}} = -0.0591\times7 = -0.414\,\text{V}$; $E_{\text{cell}} = 1.23 - (-0.414) = +1.644\,\text{V} > 0$ \checkmark
\end{solution}

\begin{problembox}
\textbf{MC.3} \hfill \textbf{[NA Level-3 Q5 / NRJ Ex-II Multi Q12]}

In the electrolysis of an aqueous \ce{NaCl} solution with inert platinum electrodes, which of the following statements are \textbf{correct}?
\begin{enumerate}[label=(\alph*)]
    \item In dilute solution, \ce{H2} is produced at the cathode.
    \item In dilute solution, \ce{O2} is produced at the anode.
    \item In concentrated solution, \ce{Cl2} is produced at the anode.
    \item Na metal is deposited at the cathode in both dilute and concentrated solutions.
\end{enumerate}
\end{problembox}
\begin{solution}
\textbf{Correct: (a), (b), (c)}

(a) True: \ce{H+} is preferentially reduced over \ce{Na+} in both dilute and concentrated. \checkmark\\
(b) True: In dilute NaCl on Pt, oxygen overpotential is not enough to divert; \ce{OH-}/\ce{H2O} is oxidized. \checkmark\\
(c) True: High [\ce{Cl-}] + high $\eta_{\ce{O2}}$ on graphite/Pt at high current density $\Rightarrow$ \ce{Cl2}. \checkmark\\
(d) False: Na metal cannot be deposited from aqueous solution (water is reduced first; Na would react violently with water anyway). \texttimes
\end{solution}

\begin{problembox}
\textbf{MC.4} \hfill \textbf{[NRJ Ex-II Multi Q16]}

Which of the following are correct statements about corrosion of iron?
\begin{enumerate}[label=(\alph*)]
    \item It is an electrochemical process involving local anodic and cathodic regions.
    \item Salt water accelerates corrosion by increasing ionic conductivity.
    \item Connecting iron to a copper plate slows corrosion of iron.
    \item Rusting requires both oxygen and water simultaneously.
\end{enumerate}
\end{problembox}
\begin{solution}
\textbf{Correct: (a), (b), (d)}

(a) True: Corrosion proceeds via local galvanic cells (e.g., grain boundaries as anodes). \checkmark\\
(b) True: NaCl ions increase the electrolyte conductivity, allowing faster ion transport between anodic and cathodic regions, greatly accelerating corrosion. \checkmark\\
(c) False: Cu is less active than Fe ($E^\circ_{\ce{Cu}} = +0.34\,\text{V}$, $E^\circ_{\ce{Fe}} = -0.44\,\text{V}$); Fe acts as anode in the galvanic couple $\Rightarrow$ Fe corrodes faster. \texttimes\\
(d) True: Both are required as established. \checkmark
\end{solution}

\begin{problembox}
\textbf{MC.5} \hfill \textbf{[PRSN Ex-II Q12 / CNG-TH Advanced]}

For the half-reaction $\ce{MnO4- + 8H+ + 5e- -> Mn^2+ + 4H2O}$ ($E^\circ = +1.51\,\text{V}$), which of the following are true?
\begin{enumerate}[label=(\alph*)]
    \item Lowering the pH increases $E$ for this half-reaction.
    \item Raising $[\ce{Mn^2+}]$ decreases $E$.
    \item This half-reaction can oxidize \ce{Cl-} to \ce{Cl2} under standard conditions.
    \item $E^\circ$ for the reverse reaction is $-1.51\,\text{V}$.
\end{enumerate}
\end{problembox}
\begin{solution}
\textbf{Correct: (b), (c), (d)}

(a) False: Nernst equation shows $E = E^\circ - \frac{0.0591}{5}\log\frac{[\ce{Mn^2+}]}{[\ce{MnO4-}][\ce{H+}]^8}$. Lowering $[\ce{H+}]$ (increasing term in denominator) \textit{decreases} $E$. \texttimes\\
(b) True: Higher $[\ce{Mn^2+}]$ increases the numerator in the log term $\Rightarrow$ $E$ decreases. \checkmark\\
(c) True: $E^\circ(\ce{Cl2/Cl-}) = +1.36\,\text{V} < 1.51\,\text{V}$, so \ce{MnO4-} can oxidize \ce{Cl-} to \ce{Cl2} (standard condition). \checkmark\\
(d) True: Reversing the half-reaction reverses the sign of $E^\circ$. \checkmark
\end{solution}
